\documentclass[conference,a4paper]{IEEEtran}
\usepackage{amsmath}
\usepackage{booktabs}
\usepackage{array}
\usepackage{url}
\usepackage[hidelinks]{hyperref}
% Generated by experiments/render_results.py; reference values are regenerated in CI.
\newcommand{\ResultSessions}{60}
\newcommand{\ResultPlays}{1080}
\newcommand{\ResultValidSessions}{60}
\newcommand{\ResultSchemaValid}{60}
\newcommand{\ResultExactRequests}{360}
\newcommand{\ResultRedactedRequests}{360}
\newcommand{\ResultHashedRequests}{360}
\newcommand{\ResultVerifiedGroundings}{360}
\newcommand{\ResultCommittedGroundings}{720}
\newcommand{\ResultStateEquivalent}{1080}
\newcommand{\ResultEventEquivalent}{1080}
\newcommand{\ResultDifferentialCases}{5000}
\newcommand{\ResultDifferentialNormalisation}{5000}
\newcommand{\ResultDifferentialState}{5000}
\newcommand{\ResultDifferentialEvents}{5000}
\newcommand{\ResultDifferentialPaths}{5000}
\newcommand{\ResultFaultClasses}{11}
\newcommand{\ResultFaults}{1560}
\newcommand{\ResultFaultsDetected}{1560}
\newcommand{\ResultSerializationVariants}{180}
\newcommand{\ResultSerializationAccepted}{180}
\newcommand{\ResultRawBytesFull}{4430}
\newcommand{\ResultRawBytesRedacted}{4797}
\newcommand{\ResultRawBytesHashed}{2259}
\newcommand{\ResultGzipBytes}{482}
\newcommand{\ResultVerificationMs}{24.73}
\newcommand{\ResultThroughput}{728}
\newcommand{\ResultTimingRepetitions}{7}
\newcommand{\ResultFaultRate}{100\%}


\title{From Prompt to State: A Domain-Semantic Reproducibility Contract for Human--LLM-Mediated Control}

\author{\IEEEauthorblockN{Tomas Laurenzo}
\IEEEauthorblockA{Department of Critical Media Practices\\
University of Colorado Boulder\\
Boulder, Colorado, USA\\
tomas@laurenzo.net}}

\begin{document}
\maketitle

\begin{abstract}
Systems that place a large language model between human intention and executable action cross two different computational regimes: stochastic language generation and deterministic state transition. Conventional chat logs preserve neither this boundary nor enough evidence to verify its consequences. This paper defines a domain-semantic reproducibility contract that separates invocation provenance from operative replay. A trace records the ordered request, inference configuration, returned text, response-to-command grounding, frozen rules, semantic events, and pre- and post-action states. Verification then checks four levels: structural validity, request reconstructability, grounding consistency, and replay-derived state equivalence. We implement the contract in BatLLM, an open-source game in which human players control agents indirectly through local language models. The implementation provides a versioned schema, three privacy modes, canonical commitments, a provider-neutral recorder, a headless verifier, and a reproducibility corpus. Across \ResultSessions{} sessions and \ResultPlays{} transitions, all traces verified for state and event equivalence. Full traces verified \ResultVerifiedGroundings{} response-to-command groundings; privacy-reduced traces retained operative replay through commitments. After every commitment derivable from retained content was recomputed, the verifier detected all \ResultFaults{} applicable controlled perturbations. The contribution is not generic record-and-replay for agents, but a contract for verifying the domain-level consequence of grounded language without assuming that model outputs can be regenerated.
\end{abstract}

\begin{IEEEkeywords}
large language models, reproducibility, provenance, record and replay, human--AI interaction, game systems
\end{IEEEkeywords}

\section{Introduction}
Large language models (LLMs) increasingly mediate actions in software environments. A human request may be expanded with system instructions and world state, submitted to a model, interpreted as a command, and executed by a deterministic program. The resulting action is often treated as if it were one operation. It is not. It crosses a probabilistic boundary and then a semantic boundary.

Even nominally deterministic decoding does not guarantee identical outputs across inference configurations. Hardware, batching, numerical precision, and execution order may alter early token probabilities and produce divergent responses~\cite{yuan2025nondeterminism}. Consequently, reproducing an LLM-mediated action by calling the model again is an unsound default. A stronger and narrower question is possible: can a system reconstruct what was sent, preserve what was returned, verify how the text was grounded, and re-execute the grounded command against the original domain semantics?

We answer this question through BatLLM~\cite{laurenzo2026batllm}, a turn-based game in which two humans specify strategies in natural language while local LLMs issue actions on their behalf. BatLLM is useful here not because games are novel agent environments---benchmarks such as AgentBoard and Cradle already use interactive software to study agent behaviour~\cite{ma2024agentboard,tan2025cradle}---but because its compact command grammar and deterministic state model expose the language-to-action boundary.

This paper contributes: (1) a formal trace model separating invocation provenance from operative reproducibility; (2) a four-level verification contract; (3) an open reference implementation with full, redacted, and hash-only privacy modes; and (4) an evaluation of replay fidelity, semantic perturbation detection, and instrumentation overhead. We make no claim that preserved prompts can regenerate identical model text. We verify the computational consequence of the text actually returned.

\section{Related Work and Scope}
W3C PROV provides interoperable concepts for entities, activities, agents, and derivation~\cite{w3c2013provo}. OpenTelemetry's generative-AI conventions further specify ordered input messages, model attributes, and output messages, while warning that recorded content may contain sensitive information~\cite{otel2026genai}. These standards establish provenance and observability vocabularies but do not define application-domain transition semantics.

Recent systems address more of the execution chain. PROV-AGENT integrates prompts, responses, decisions, and downstream workflow provenance~\cite{souza2025provagent}. AgentRR records trajectories and abstracts them into reusable experiences~\cite{feng2025agentrr}. Prompt Scene Investigation introduces chain-of-custody and controlled behavioural replay for forensic analysis~\cite{ledjaki2026psi}. Mudasiru's \emph{agrepl} captures external interactions at the transport layer and replays them without outbound access~\cite{mudasiru2026agrepl}. Paduraru \emph{et al.} define Message--Action Traces, executable contracts, deterministic replay, fault injection, and action mediation for multi-agent orchestration~\cite{paduraru2026assurance}.

These antecedents preclude any defensible claim that BatLLM introduces provenance, contracts, or record-and-replay for LLM agents. Table~\ref{tab:scope} locates the narrower contribution. Transport replay reproduces observed I/O; workflow provenance represents lineage; behavioural replay studies model outputs; orchestration contracts constrain multi-agent execution. The present work verifies whether a retained linguistic output was grounded into the recorded command and whether that command, executed by frozen domain semantics, derives the recorded state and semantic events.

\begin{table*}[t]
\caption{Primary verification object in closely related approaches. A dash does not indicate a deficiency; it marks a different research objective.}
\label{tab:scope}
\centering
\small
\begin{tabular}{@{}p{0.18\textwidth}p{0.22\textwidth}p{0.25\textwidth}p{0.27\textwidth}@{}}
\toprule
Approach & Primary recorded object & Replay or assurance target & Distinction from this contract \\
\midrule
W3C PROV / OpenTelemetry~\cite{w3c2013provo,otel2026genai} & Provenance graph; GenAI invocation telemetry & Lineage and observability & No application-specific grounding or transition-equivalence rule \\
PROV-AGENT~\cite{souza2025provagent} & Agent interactions in workflow provenance & End-to-end provenance queries and reliability analysis & Relates decisions and outcomes; does not re-derive a bounded domain state \\
AgentRR~\cite{feng2025agentrr} & Successful interaction trajectory and abstracted experience & Reuse of prior experience in later tasks & Replay guides new execution rather than verifying the original transition \\
PSI~\cite{ledjaki2026psi} & Prompt--response evidence with chain of custody & Controlled model-behaviour replay and investigation & Re-enacts model behaviour; does not verify text-to-command domain semantics \\
\emph{agrepl}~\cite{mudasiru2026agrepl} & External transitions & Isolated deterministic I/O replay & Reproduces dependencies below the application semantic boundary \\
MAT contracts~\cite{paduraru2026assurance} & Message--action orchestration traces & Contract verdicts, deterministic replay, fault containment & General orchestration assurance; this paper specifies post-grounding state/event equivalence \\
BatLLM contract & Invocation, response, grounded command, rules, events, states & Exact request reconstruction where retained; grounding and domain-state verification & Explicitly couples stochastic language evidence to a deterministic transition function \\
\bottomrule
\end{tabular}
\end{table*}

q���q�^